\documentclass[12pt,a4paper]{article}
\usepackage{cmap} % Makes the PDF copiable. See http://tex.stackexchange.com/a/64198/25761
\usepackage[T1]{fontenc}
\usepackage[brazil]{babel}
\usepackage[utf8]{inputenc}
\usepackage{amsmath}
\usepackage{amsfonts}
\usepackage{amssymb}
\usepackage{amsthm}
\usepackage{textcomp} % \degree
\usepackage{gensymb} % \degree
\usepackage[usenames,svgnames,dvipsnames]{xcolor}
\usepackage{hyperref}
\usepackage{graphicx}
\usepackage[margin=2cm]{geometry}

\hypersetup{
    colorlinks = true,
    allcolors = {blue}
}
\usepackage{cancel}

% TODO: Consider using exsheets
% http://linorg.usp.br/CTAN/macros/latex/contrib/exsheets/exsheets_en.pdf
%
% http://ctan.org/tex-archive/macros/latex/contrib/exercise/
% Options: answerdelayed,lastexercise,noanswer
\usepackage[answerdelayed,lastexercise]{exercise}

\addto\captionsbrazil{%
\def\listexercisename{Lista de exerc\'icios}%
\def\ExerciseName{Exerc\'icio}%
\def\AnswerName{Solu\c{c}\~ao do exerc\'icio}%
\def\ExerciseListName{Ex.}%
\def\AnswerListName{Solu\c{c}\~ao}%
\def\ExePartName{Parte}%
\def\ArticleOf{de\ }%
}

\renewcommand{\ExerciseHeaderTitle}{(\ExerciseTitle)\ }
\renewcommand{\ExerciseListHeader}{%\ExerciseHeaderDifficulty%
\textbf{%\ExerciseListName\
\ExerciseHeaderNB.\ %
%\ --- \
\ExerciseHeaderTitle}%
%\ExerciseHeaderOrigin
\ignorespaces}
\renewcommand{\AnswerListHeader}{\textbf{\ExerciseHeaderNB.\ (\AnswerListName)\ }}

\newcommand*\sen{\operatorname{sen}}
\newcommand*\cotg{\operatorname{cotg}}
\newcommand*\cotgh{\operatorname{cotgh}}
\newcommand*\dom[1]{\operatorname{Dom}\left(#1\right)}

\newcommand*\R{\mathbb{R}}

\renewcommand{\theenumi}{\alph{enumi}}
\renewcommand\labelenumi{(\theenumi) }

\newcommand*\tipo{Prova I}
\newcommand*\turma{TADS121-01C}
\newcommand*\disciplina{CDI0001}
\newcommand*\eu{Helder G. G. de Lima}
\newcommand*\data{09/09/2016}

\author{\eu}
\title{\tipo - \disciplina}
\date{\data}

\begin{document}
\thispagestyle{empty}
\newgeometry{margin=2cm,bottom=0.5cm}
\begin{center}
\includegraphics[width=9.0cm]{marca} \\
\textbf{\tipo\ (\disciplina / \turma)} \\
Prof. \eu\footnote{
Este é um material de acesso livre distribuído sob os termos da licença \href{https://creativecommons.org/licenses/by-sa/4.0/deed.pt_BR}{Creative Commons BY-SA 4.0}.}
\end{center}

\noindent Nome do(a) aluno(a): \underline{\hspace{9,7cm}} Data: \underline{\data}

%\section*{Instruções}
\begin{center}\fbox{
\begin{minipage}{14cm}

{\footnotesize
\begin{itemize}
\renewcommand{\theenumi}{\Roman{enumi}}
\item Identifique-se em todas as folhas.
\item Mantenha o celular e os demais equipamentos eletrônicos desligados durante a prova.
\end{itemize}
}

\end{minipage}
}
\end{center}


%\section*{Questões}
\begin{ExerciseList}
\Exercise[title={2,5}] Sejam $p$ e $q$ as funções definidas pelas regras $p(x) = \dfrac{1}{5+x}$ e $q(x) = \dfrac{15}{1-x^2}$ (nos pontos $x \in \R$ em que estas expressões fazem sentido).
\begin{enumerate}
\item \textbf{(0,5)} Obtenha o domínio das funções $p$ e $q$.
\item \textbf{(0,5)} Levando em conta o item anterior, qual é o domínio da função composta $p \circ q$?
\item \textbf{(0,5)} Verifique se $p\circ q$ é par ou ímpar, ou se não tem nenhum destes tipos de simetrias.
\item \textbf{(1,0)} Para cada ponto $a > 0$ que não pertencer ao domínio de $p\circ q$, calcule (se existirem):
\[
\lim_{x \to a^-} (p\circ q)(x)
\quad \text{e} \quad
\lim_{x \to a^+} (p\circ q)(x).
\]
\end{enumerate}
\Answer
\begin{enumerate}
\item As regras de $p$ e $q$ só não fazem sentido se os denominadores forem zero, ou seja, se $5+x = 0$, no caso de $p$, e se $1-x^2 = 0$, no caso de $q$. Então $\dom{p} = \R \setminus \{ -5\}$ e $\dom{q} = \R \setminus \{ -1, +1 \}$.
\item Para que $x \in \dom{p \circ q}$ é preciso que $x \in \dom{q}$ e que além disso $q(x) \in \dom{p}$. Em outras palavras, deve ocorrer que $x \neq \pm 1$ e que $q(x) \neq -5$. Mas para $x \neq \pm 1$, tem-se
\[
\dfrac{15}{1-x^2} = -5
\Leftrightarrow 15 = -5( 1-x^2 )
\Leftrightarrow -3 = 1-x^2
\Leftrightarrow x^2 = 4
\Leftrightarrow x = \pm 2.
\]
Portanto,
\[
\dom{p \circ q} = \R \setminus \{-2, -1, 1, 2\} = (-\infty, -2) \cup (-2,-1) \cup (-1,1) \cup (1,2) \cup (2, +\infty).
\]
\item Para $x \in \dom{p \circ q}$ tem-se
\[
(p \circ q) (x)
= p(q(x))
= \dfrac{1}{5+q(x)}
= \dfrac{1}{5+\dfrac{15}{1-x^2}}
= \dfrac{1}{\dfrac{20-5x^2}{1-x^2}}
= \dfrac{1-x^2}{20-5x^2}.
\]
Então
\[
(p \circ q) (-x)
= \dfrac{1-(-x)^2}{20-5(-x)^2}
= \dfrac{1-x^2}{20-5x^2}
= (p \circ q) (x),
\]
ou seja, $p \circ q$ é uma função par.
\item Pelo que foi obtido anteriormente, os únicos valores de $a>0$ que não pertencem ao domínio de $p \circ q$ são $a = 1$ e $a = 2$. Como $(p \circ q)(x) = \dfrac{1-x^2}{20-5x^2} = \dfrac{(1-x)(1+x)}{5(2+x)(2-x)}$, tem-se:
\[
\lim_{x \to 1^-} \dfrac{(1-x)(1+x)}{5(2+x)(2-x)}
= \dfrac{(1-1)(1+1)}{5(2+1)(2-1)}
= 0
=
\lim_{x \to 1^+} \dfrac{(1-x)(1+x)}{5(2+x)(2-x)}.
\]
Já em $a=2$, ocorre o seguinte:
\[
\lim_{x \to 2^-} \dfrac{(1-x)(1+x)}{5(2+x)(2-x)}
= \lim_{x \to 2^-} \cancelto{\frac{-3}{20}}{ \dfrac{(1-x)(1+x)}{5(2+x)} }
\quad
\cdot \cancelto{ +\infty }{\frac{1}{2-x}}
= - \infty
\]
e
\[
\lim_{x \to 2^+} \dfrac{(1-x)(1+x)}{5(2+x)(2-x)}
= \lim_{x \to 2^+} \cancelto{\frac{-3}{20}}{ \dfrac{(1-x)(1+x)}{5(2+x)} }
\quad
\cdot \cancelto{ -\infty }{\frac{1}{2-x}}
= +\infty.
\]
\end{enumerate}

\Exercise[title={2,5}] Sejam $f(x) = \dfrac{2x^2-x-1}{x-1}$ e $\displaystyle L = \lim_{x \to 1} f(x)$.
\begin{enumerate}
\item \textbf{(1,0)} Calcule o valor de $L$.
\item \textbf{(1,5)} Verifique que para o valor de $L$ obtido, pode-se encontrar para cada $\varepsilon > 0$ algum valor de $\delta > 0$ tal que
$|f(x) - L| < \varepsilon$ sempre que $0 < |x-1| < \delta$.
\end{enumerate}
\Answer
\begin{enumerate}
\item Se $x \neq 1$, então $\dfrac{2x^2-x-1}{x-1} = \dfrac{(2x+1)(x-1)}{x-1} = 2x+1$. Portanto:
\[
L = \lim_{x \to 1} f(x) = \lim_{x \to 1} = \lim_{x \to 1} 2x+1 = 3.
\]

\textbf{Solução alternativa}:
\[
L = \lim_{x \to 1} \dfrac{2x^2-x-1}{x-1}
  = \lim_{u \to 0} \dfrac{2(u+1)^2-(u+1)-1}{u}
  = \lim_{u \to 0} \dfrac{2u^2 + 3u}{u}
  = \lim_{u \to 0} 2u+3 = 3.
\]

\item Pelo item anterior, tem-se:
\[
|f(x) - L| < \varepsilon
\Leftrightarrow \left|(2x+1) - 3 \right| < \varepsilon
\Leftrightarrow |2x-2| < \varepsilon
\Leftrightarrow 2|x-1| < \varepsilon
\Leftrightarrow |x-1| < \frac{\varepsilon}{2}.
\]
Logo, dado $\varepsilon > 0$, basta considerar $\delta = \frac{ \varepsilon}{2}$, e ele terá a propriedade desejada.

\end{enumerate}

\Exercise[title={2,0}] Determine os valores de $a$ e $b$ para que a função a seguir seja contínua em 0:
\[
g(x) =
\begin{cases}
a + \frac{ \tan(x) }{ x },\ & \text{se}\ x > 0,\\
b,\ &\text{se}\ x = 0\\
\frac{ \sen(3x) }{ \sen(x) },\ &\text{se}\ x < 0
\end{cases}
\]
\Answer Para que $g$ seja contínua no ponto indicado, é preciso que $\lim_{x \to 0} g(x) = g(0)$, e para que este limite exista, é necessário que os limites laterais coincidam. Mas
\begin{align*}
  \lim_{x \to 0^-} g(x)
& = \lim_{x \to 0^-} \frac{ \sen(3x) }{ \sen(x) }
= \lim_{x \to 0^-} \frac{ \frac{ \sen(3x) }{x} }{ \frac{ \sen(x) }{x} }
= \frac{ \lim_{x \to 0^-} \frac{ \sen(3x) }{x} }
       { \lim_{x \to 0^-} \frac{ \sen( x) }{x} }\\
&= \frac{ \lim_{x \to 0^-} 3 \frac{ \sen(3x) }{3x} }
       { \lim_{x \to 0^-} \frac{ \sen( x) }{x} }
= \frac{ 3 \lim_{u \to 0^-} \frac{ \sen(u) }{u} }
       { \lim_{x \to 0^-} \frac{ \sen( x) }{x} }
= 3.
\end{align*}
Logo, é preciso que $b = g(0) = 3$. Por outro lado,
\begin{align*}
    \lim_{x \to 0^+} g(x)
& = \lim_{x \to 0^+} a + \frac{ \tan(x) }{ x }
  = a + \lim_{x \to 0^+} \frac{ \frac{ \sen(x) }{ \cos{x} } }{ x }\\
& = a + \lim_{x \to 0^+} \frac{ \sen(x) }{ x \cos{x} }
  = a + \lim_{x \to 0^+} \frac{ \sen(x) }{ x }\frac{1}{ \cos{x} }
  = a + 1.
\end{align*}
Portanto, para que o limite exista também é necessário que $a+1 = 3$, ou seja, $a=2$.

\Exercise[title={3,0}] Calcule 2 dos seguintes limites, se existirem (ou explique porque não existem):
\begin{enumerate}
\item (1,5) $\lim\limits_{x\to 0} 2x \cotg(5x)$
\item (1,5) $\lim\limits_{x\to 1/3} \dfrac{\ln(x) - \ln(1/3)}{x-1/3}$
\item (1,5) $\lim\limits_{x\to +\infty} \cotgh(x)$, sendo $\cotgh(x) = \dfrac{e^x+e^{-x}}{e^x-e^{-x}}$
\end{enumerate}
\Answer
\begin{enumerate}
\item
\begin{align*}
  \lim\limits_{x\to 0} 2x \cotg(5x)
&= \lim\limits_{x\to 0} 2x \frac{ \cos(5x) }{ \sen(5x) }
= \lim\limits_{u\to 0} 2\frac{u}{5} \frac{ \cos(u) }{ \sen(u) }
= \frac{2}{5} \lim\limits_{u\to 0} \frac{ u }{ \sen(u) }\cos(u)\\
&= \frac{2}{5} \lim\limits_{u\to 0} \frac{ 1 }{  \frac{ \sen(u) }{ u } }\cos(u)
= \frac{2}{5}.
\end{align*}

\item
\begin{align*}
  \lim\limits_{x\to 1/3} \dfrac{\ln(x) - \ln(1/3)}{x-1/3}
& = \lim\limits_{u\to 0} \dfrac{ \ln( u + 1/3 ) - \ln(1/3) }{u}
= \lim\limits_{u\to 0} \dfrac{ \ln\left( \frac{ u + 1/3 }{ 1/3 } \right)}{u}
= \lim\limits_{u\to 0} \frac{1}{u} \ln\left( 3u + 1 \right) \\
& = \lim\limits_{u\to 0}  \ln\left( 3u + 1 \right)^{\frac{1}{u}}
 = \ln\left( \lim\limits_{u\to 0} \left( 3u + 1 \right)^{\frac{1}{u}} \right)
 = \ln\left( \lim\limits_{t\to 0} \left( 1 + t \right)^{\frac{3}{t}} \right)\\
& = \ln\left( \lim\limits_{t\to 0} \left[ \left( 1 + t \right)^{\frac{1}{t}} \right]^3 \right)
  = \ln\left( \left[ \lim\limits_{t\to 0} \left( 1 + t \right)^{\frac{1}{t}} \right]^3 \right)
  = \ln(e^3)
  = 3.
\end{align*}

\item
\begin{align*}
\lim\limits_{x\to +\infty} \cotgh(x)
& = \lim\limits_{x\to +\infty} \dfrac{e^x+e^{-x}}{e^x-e^{-x}}
= \lim\limits_{x\to +\infty} \dfrac{e^x+\frac{1}{e^x}}{e^x-\frac{1}{e^x}}
= \lim\limits_{w\to +\infty} \dfrac{w+\frac{1}{w}}{w-\frac{1}{w}}
= \dfrac{\lim\limits_{w\to +\infty}  w+\frac{1}{w}}{\lim\limits_{w\to +\infty}  w-\frac{1}{w}}
= \frac{1}{1}
= 1.
\end{align*}

\end{enumerate}
\end{ExerciseList}

\begin{center}
BOA PROVA!
\end{center}

\newpage
\restoregeometry
\section*{Respostas}
\shipoutAnswer
\end{document}
